\section{Categorical duality: the Chu construction}\label{sec:chu}

Every duality encountered in this paper pairs a ``forward'' object with
a ``backward'' object through a scalar-valued coupling:
cost-to-go $P$ with covariance $\Sigma$ via the trace pairing $\tr(P\Sigma)$
(\S\ref{sec:riccati});
value function $V$ with the linearized solution
$\psi = e^{-V/(2\varepsilon)}$ via Hopf--Cole (\S\ref{sec:mppi});
a function $f$ on an LCA group with its Fourier transform $\hat{f}$
via the character pairing (\S\ref{sec:cycle});
the causal Hardy space $\Htwo$ with its anti-causal complement
$(\Htwo)^\perp$ (\S\ref{sec:hardy});
representation spaces $V_\rho$ with their contragredients $V_\rho^*$
(item~(v) below);
velocities $\xi \in \mathfrak{g}$ with momenta $\mu \in \mathfrak{g}^*$
via the Legendre transform (\S\ref{sec:articulated}).
In each case, the morphisms---Fourier transform, Hopf--Cole substitution,
Legendre transform---act forward on one component and contravariantly on
the dual.

These are not independent observations.
They are instances of a single categorical construction due to
Barr~\cite{Barr1979,BarrOmatt}, whose antecedents include the
``Mackey pairs'' of~\cite{Mackey1945}.

\subsection*{The Chu construction}

\begin{definition}[Chu space {\cite[Definition~2.1]{BarrOmatt}}]
\label{def:chu-space}
Let $(\mathcal{V}, \otimes, I, \multimap)$ be a symmetric monoidal
closed category with pullbacks and let
$\bot \in \mathrm{Ob}(\mathcal{V})$ be a fixed object (the
\emph{dualizing object}).
The category $\Chu(\mathcal{V}, \bot)$ has:
\begin{enumerate}[label=(\roman*),nosep]
\item \emph{Objects:} triples $(A, r, X)$ with
  $A, X \in \mathrm{Ob}(\mathcal{V})$ and
  $r \colon A \otimes X \to \bot$ (the \emph{pairing}).
\item \emph{Morphisms} $(A, r, X) \to (B, s, Y)$: pairs
  $(f \colon A \to B,\; g \colon Y \to X)$ such that
  \[
    s \circ (f \otimes \mathrm{id}_Y)
    = r \circ (\mathrm{id}_A \otimes g)
    \colon A \otimes Y \to \bot.
  \]
  The first component $f$ goes \emph{forward}; the second $g$ goes
  \emph{backward}.
\end{enumerate}
The dual of $(A, r, X)$ is $(X, r^\sigma, A)$, where $r^\sigma$ is
the transpose $X \otimes A \xrightarrow{\sim} A \otimes X
\xrightarrow{r} \bot$.
\end{definition}

\begin{theorem}[{Barr \cite[Theorem~4.1]{BarrOmatt}}]
\label{thm:chu-star-auto}
If $\mathcal{V}$ is a complete symmetric monoidal closed category with
pullbacks and $\bot \in \mathrm{Ob}(\mathcal{V})$, then
$\Chu(\mathcal{V}, \bot)$ is a $*$-autonomous category: it is
symmetric monoidal closed with dualizing object
$(\bot, \mathrm{ev}, I)$, and the canonical map
$\mathbf{A} \to \mathbf{A}^{**}$ is a natural isomorphism.
\end{theorem}

A $*$-autonomous category is a symmetric monoidal closed category
$(\mathcal{C}, \otimes, I, \multimap)$ equipped with a dualizing object
$\bot$ such that $A \mapsto (A \multimap \bot) \multimap \bot$ is
naturally isomorphic to the identity.
The relevance to linear logic is that $*$-autonomous categories are
the categorical models of classical multiplicative linear
logic~\cite{Seely1989,Girard}.

\begin{remark}[Separated and extensional]
\label{rmk:sep-ext}
A Chu object $(A, r, X)$ is \emph{separated} if
$r(a, -) = r(a', -)$ implies $a = a'$
(the pairing separates points of $A$) and \emph{extensional} if
$r(-, x) = r(-, x')$ implies $x = x'$.
The full subcategory of separated-extensional objects is denoted
$\mathrm{chu}(\mathcal{V}, \bot)$; it is still
$*$-autonomous~\cite[Theorem~4.1(7)]{BarrOmatt}.
This refinement eliminates degeneracies: the Pontryagin duality
$G \leftrightarrow \widehat{G}$ for LCA groups, for instance, lives in
$\mathrm{chu}(\mathsf{Ab}, \bT)$, where only non-degenerate pairings
$G \times \widehat{G} \to \bT$ are
admitted~\cite[Example~5.3]{BarrOmatt}.
\end{remark}

\subsection*{The eight Chu objects of optimal cycle theory}

Every duality in Parts~I and~II defines a separated Chu object.
We record them as triples
$(A,\; r \colon A \otimes X \to \bot,\; X)$ with specified
$\bot$.

\begin{proposition}\label{prop:eight-chu}
The following are separated Chu objects in the indicated ambient
categories:
\begin{enumerate}[label=\textup{(\roman*)},nosep]
\item \emph{Control--estimation Riccati}
  \textup{(\S\ref{sec:riccati}):}
  $A = \mathrm{Sym}^+_n$ (cost-to-go matrices $P$),
  $X = \mathrm{Sym}^+_n$ (covariance matrices $\Sigma$),
  $r(P, \Sigma) = \tr(P\Sigma)$,
  $\bot = \R$.
  The control Riccati and the Kalman--Bucy Riccati solve the same
  equation with $P$ and $\Sigma$ exchanged.

\item \emph{Hopf--Cole linearization}
  \textup{(\S\ref{sec:mppi}):}
  $A = C^\infty(G)$ (value functions $V$),
  $X = C^\infty_{>0}(G)$ (linear solutions
  $\psi = e^{-V/(2\varepsilon)}$),
  $r(V, \psi) = \int_G \psi\, e^{V/(2\varepsilon)} d\mu_G$,
  $\bot = \R$.
  The substitution $V \leftrightarrow \psi$ is an isomorphism of
  Chu objects; the nonlinear HJB becomes the linear heat equation.

\item \emph{Pontryagin--Fourier duality}
  \textup{(\cref{thm:pontryagin}):}
  $A = G$ (LCA group),
  $X = \widehat{G}$ (Pontryagin dual),
  $r(x, \chi) = \chi(x) \in \bT$,
  $\bot = \bT$.
  The evaluation pairing $G \times \widehat{G} \to \bT$ is exactly the
  Chu pairing; the Pontryagin isomorphism $G \cong
  \widehat{\widehat{G}}$ is the $*$-autonomy $A \cong A^{**}$.

\item \emph{Hardy space decomposition}
  \textup{(\S\ref{sec:hardy}):}
  $A = \Htwo(\bT)$ (causal),
  $X = (\Htwo)^\perp$ (anti-causal),
  $r(f_+, f_-) = \int_\bT f_+ \bar{f}_-\, d\theta$,
  $\bot = \mathbb{C}$.
  The Szeg\H{o} projection
  $L^2(\bT) = \Htwo \oplus (\Htwo)^\perp$
  is the Chu decomposition into forward-causal and backward-causal
  components.

\item \emph{Peter--Weyl blocks}
  \textup{(\S\ref{sec:compact}):}
  $A = V_\rho$ (representation space),
  $X = V_\rho^*$ (contragredient),
  $r(v, v^*) = v^*(v)$,
  $\bot = \mathbb{C}$.
  The matrix coefficient
  $\rho_{ij}(g) = e_j^*(\ \rho(g)\, e_i)$ is a morphism of Chu
  objects; the Schur orthogonality relations express the inner product
  in $\mathrm{chu}(\mathsf{Vect}, \mathbb{C})$.

\item \emph{Lie--Poisson duality}
  \textup{(\S\ref{sec:articulated}):}
  $A = \mathfrak{g}$ (velocities),
  $X = \mathfrak{g}^*$ (momenta),
  $r(\xi, \mu) = \langle \mu, \xi \rangle$,
  $\bot = \R$.
  The Legendre transform
  $\mathbb{FL} \colon \mathfrak{g} \to \mathfrak{g}^*$ and its inverse
  constitute a Chu isomorphism.
  The Euler--Poincar\'e equation on $\mathfrak{g}$ and the Lie--Poisson
  equation on $\mathfrak{g}^*$ are the forward and backward faces.

\item \emph{Barrier--Boltzmann coupling}
  \textup{(\S\ref{sec:articulated}):}
  $A = C(Q)$ (barrier potentials $-\varepsilon\log h$),
  $X = \cP(Q)$ (Boltzmann measures
  $w \propto e^{-J/\varepsilon}$),
  $r(V, w) = \int_Q V\, dw = \mathbb{E}_w[V]$,
  $\bot = \R$.
  The nonuple parameter $\varepsilon$ of \S\ref{sec:articulated}
  scales the pairing: as $\varepsilon \to 0$, the Boltzmann measure
  $w$ concentrates on the support of the barrier constraint.

\item \emph{Action--state duality}
  \textup{(\S\ref{sec:articulated}, eq.~\eqref{eq:action-state}):}
  $A = \R$ (log-contact modes $\beta_k$, the Lie algebra of
  generators---lift-off, touch-down),
  $X = (0,1)$ (contact modes $\sigma_k$, the Lie group of
  configurations),
  $r(\beta, \sigma) = \beta\,\sigma$ (the tangent-value pairing),
  $\bot = \R$.
  The sigmoid $\sigma = (1 + e^{-\beta})^{-1}$ is the exponential map;
  its inverse $\beta = \log(\sigma/(1{-}\sigma))$ is the logit.
  The forward map $\beta \to x(T)$ (given $x(0)$) is a function;
  the inverse $x(T) \to \beta^*$ requires the variational principle
  \eqref{eq:least-action}---the Chu adjunction condition.
\end{enumerate}
\end{proposition}

\begin{proof}
In each case the verification is the adjunction condition of
\cref{def:chu-space}(ii): a morphism $(f, g)$ satisfies
$s \circ (f \otimes \mathrm{id}) = r \circ (\mathrm{id} \otimes g)$
if and only if the map $f$ and its transpose $g$ are related by the
pairing.
For~(i) this is $\tr(KP\Sigma) = \tr(P K^\top \Sigma)$;
for~(iii) this is the functoriality of the Fourier transform
(\cref{rmk:functoriality});
for~(vi) this is the definition of the Legendre--Fenchel transform.
The separation property follows from the non-degeneracy of each pairing
(e.g., $\tr(PX) = 0$ for all $X \succeq 0$ implies $P = 0$).
The remaining cases are analogous.
\end{proof}

\begin{remark}[Morphisms between Chu objects]
The key transformations of the paper are Chu morphisms:
\begin{enumerate}[label=(\alph*),nosep]
\item The \emph{Fourier transform}
  $\cF \colon (G, r, \widehat{G}) \to (G, r, \widehat{G})$ acts
  forward on functions and backward (contravariantly) on characters.
\item The \emph{Hopf--Cole substitution}
  $(V, r, \psi) \mapsto (-2\varepsilon\log\psi, r', \psi)$ is a
  Chu isomorphism exchanging the nonlinear and linear pictures.
\item The \emph{Legendre transform}
  $\mathbb{FL} \colon (\mathfrak{g}, r, \mathfrak{g}^*) \to
  (\mathfrak{g}^*, r^\sigma, \mathfrak{g})$
  swaps $A$ and $X$---it is the Chu dual map.
\end{enumerate}
The direction reversal in the second component---the defining feature
of the Chu construction---is the source of every ``duality'' in the
paper.
\end{remark}

\begin{remark}[Harmonic predictive control]
\label{rmk:hpc}
The same author~\cite{PrattHPC} independently applied the exponential
bridge to PID control.
Pratt's \emph{harmonic predictive control} (HPC) models the
trajectory as $\varepsilon(t) = e^{\beta(t)} + \delta$, where
$\beta(t)$ is a complex polynomial:
\begin{itemize}[nosep]
\item \emph{Limit avoidance} (sinusoidal): $\beta' = i\omega$,
  giving oscillation---the Fourier/character side of the cycle;
\item \emph{Setpoint seeking} (Gaussian): $\beta = -\frac{1}{2}
  \bigl(\frac{t-\mu}{\sigma}\bigr)^2$
  (real quadratic), giving heat-kernel decay---the diffusion side.
\end{itemize}
The controlled variable is the third derivative (jerk)
$j = \varepsilon'''$, computed from $(q, v, a)$ via the differential
equation $\varepsilon''' = (3\beta'' + \beta'^2)\varepsilon'$,
with no tuning parameters: the control law is \emph{ordained by the
exponential structure}.

The parallel with \S\ref{sec:articulated} is precise.
The ``proximity'' $p_\ell = v/(q - \ell)$ diverges as $q \to \ell$,
exactly as our barrier force scale $\sigma_k = \varepsilon/h_k$
diverges as $h_k \to 0$; both are log-barrier gradients.
The sinusoidal and Gaussian regimes correspond to the two sides of
the optimal cycle (\S\ref{sec:cycle}), and their $C^2$-continuous
handoff is the smooth contact transition mediated by the barrier.
Our Riccati inner solve (\S\ref{sec:articulated}) replaces
hand-tuned PID gains $(K_P, K_I, K_D)$ with optimal gains $K(t)$
derived from the same exponential bridge---the same structural move
as HPC, but via the matrix Riccati equation rather than the scalar
differential equation for $\beta$.

That the Chu construction (\cref{def:chu-space}) and the
exponential control law emerge from the same author is not a
coincidence: both are consequences of the duality
$\beta \leftrightarrow e^\beta$ between the log-domain (additive,
linear, ``tuning-free'') and the physical domain (multiplicative,
nonlinear, constrained).

The log-space B-spline parameterization \eqref{eq:log-contact} makes
this parallel precise: $\beta_k(t)$ is a periodic polynomial (B-spline)
in log space, exactly as Pratt's $\beta(t)$; and
$\sigma_k = \mathrm{sigmoid}(\beta_k) = 1/(1 + e^{-\beta_k})$ is the
smooth contact mode, exactly as Pratt's
$\varepsilon = e^\beta + \delta$.
The MPPI outer loop optimizes the B-spline coefficients
$c_{k,j} \in \R$---the ``tuning-free'' log-domain
representation---while the Riccati inner solve computes optimal gains
from the resulting $\sigma_k(t)$---the physical-domain trajectory.
The MPPI outer loop, viewed as a path integral over log-space contact
schedules (\S\ref{sec:articulated}), makes this Chu duality
computational: the forward map $\beta \to x(T)$ evaluates the Chu
pairing, while the Boltzmann-weighted inverse (path-integral mean)
is the adjunction condition.
\end{remark}

\subsection*{Linear logic interpretation}

A $*$-autonomous category carries two distinct monoidal products:
the tensor $\otimes$ and the \emph{par}
$A \parr B \coloneqq (A^* \otimes B^*)^*$, corresponding to the two
multiplicative connectives of Girard's linear logic~\cite{Girard}.
In a resource-sensitive reading:
\begin{itemize}[nosep]
\item $A \otimes B$ means ``both $A$ and $B$ are available
  simultaneously''---in control, the simultaneous availability of
  state $x$ and co-state $p$;
\item $A \parr B$ means ``$A$ or $B$, opponent's choice''---the dual
  structure where either the controller or the estimator claims the
  resource.
\end{itemize}

The exponential modality $!$ of linear logic converts a \emph{linear}
resource (use exactly once) into a \emph{classical} one (use
arbitrarily many times).
Two structures in the paper realize this conversion:

\medskip\noindent\emph{(a) The barrier method as $!$.}
The hard constraint $h_k \geq 0$ is a linear resource: the contact
force $\lambda_k$ is available in stance and consumed in swing---use
once per stride.
The log-barrier $-\varepsilon\log h_k$ converts this into a smooth
potential that can be evaluated (``used'') at every point in
configuration space.
As $\varepsilon \to 0$, the smooth potential sharpens to the hard
constraint---the $!$ ``forgets'' and the resource becomes linear
again.
The dual ascent on $\varepsilon$ (\S\ref{sec:cycle}) is the
\emph{dereliction} map $!A \to A$: each iteration tightens the
barrier, consuming one layer of the exponential modality.

\medskip\noindent\emph{(b) Hopf--Cole as $!$.}
The substitution $\psi = e^{-V/(2\varepsilon)}$ converts the
multiplicative (nonlinear) HJB into the additive (linear) heat
equation.
This is literally the exponential applied to the value function: it
converts a nonlinear/multiplicative object into a
linear/additive one.
The inverse $V = -2\varepsilon\log\psi$ is the co-exponential;
the Feynman--Kac representation (\S\ref{sec:mppi}) provides the
semantic content.

\begin{remark}[The $\varepsilon$-dereliction path]
\label{rmk:dereliction}
In the dual ascent $\varepsilon \to 0$:
\begin{enumerate}[nosep]
\item the Boltzmann distribution
  $w \propto e^{-J/\varepsilon}$ concentrates on the optimal
  trajectory (the $!$ resource is consumed);
\item the barrier $-\varepsilon\log h$ sharpens to complementarity
  ($!A \to A$: the smooth constraint becomes hard);
\item the Hopf--Cole linearization ceases to be useful
  ($\psi \to \delta$: the Feynman--Kac measure degenerates to a point
  mass).
\end{enumerate}
The $\varepsilon \to 0$ limit is the passage from the ``classical''
(resource-abundant, smooth, linear) regime to the ``linear-logic''
(resource-scarce, hard, nonlinear) regime.
The holonomy $H(f)$ (\cref{def:cycle}) measures the cost of this
passage.
\end{remark}

\subsection*{Tropical Chu spaces and the Legendre--Fenchel transform}

The Maslov dequantization~\cite{Maslov} sends the semiring
$(\R, +, \times)$ to the tropical semiring
$(\R_{\max}, \max, +)$ via the logarithmic limit
\[
\lim_{\varepsilon \to 0^+} \varepsilon \log(e^{a/\varepsilon} +
e^{b/\varepsilon}) = \max(a,b).
\]
Under this deformation, the HJB equation becomes \emph{linear over
$\R_{\max}$}: the $\min$ in the Bellman recursion is the tropical
addition, and cost accumulation along trajectories is tropical
multiplication.
The Hopf--Cole substitution (\S\ref{sec:mppi}) is precisely this
dequantization: at finite $\varepsilon$, the system lives in the
standard semiring (the heat equation); at $\varepsilon = 0$, it
passes to the tropical semiring (the Hamilton--Jacobi equation).

\begin{proposition}[Tropical Chu spaces]
\label{prop:tropical-chu}
Write $\R_{\max}$ for the tropical semiring
$(\R \cup \{-\infty\}, \max, +)$.
In $\Chu(\mathsf{Set}, \R_{\max})$, objects are pairs $(X, Y)$
with pairing $\varphi \colon X \times Y \to \R_{\max}$.
\begin{enumerate}[label=\textup{(\roman*)},nosep]
\item The Legendre--Fenchel transform
  $f^*(y) = \sup_x\, [\langle x, y \rangle - f(x)]$ arises as the
  unit of the \emph{nucleus monad} on the enriched profunctor
  $\R^{\op}_{\max} \times \R_{\max} \to \R_{\max}$
  \cite{Willerton2015}.
  The nucleus is the separated-extensional core of the Chu
  construction: a convex function $f$ and its conjugate $f^*$ form a
  separated Chu object $(f, r, f^*)$ with
  $r(x, y) = \langle x, y \rangle$ and
  $\bot = \R_{\max}$.
\item The Bellman value function $V(x) = \min_u \int L\, dt + \Phi$
  and the Lagrangian $L(x, u)$ form a Chu pair in
  $\Chu(\mathsf{Set}, \R_{\max})$ via the tropical inner product
  $\langle V, L \rangle_{\mathrm{trop}} = \max_x\,[V(x) + L(x)]$.
\end{enumerate}
\end{proposition}

Pratt~\cite{Pratt1995} observed that the dualizing object $K$ in
$\Chu(\mathsf{Set}, K)$ determines the ``physics'':
$K = \{0,1\}$ gives classical logic,
$K = [0,1]$ fuzzy/probabilistic systems,
$K = \mathbb{C}$ quantum mechanics.
The extension $K = \R_{\max}$ gives tropical/optimal-control
``physics'': ``probabilities'' become costs, ``expectations''
become optimizations.

\begin{remark}[The $\varepsilon$-interpolation of Chu categories]
\label{rmk:eps-interpolation}
At finite $\varepsilon > 0$, the dualities of this paper live in
$\Chu(\mathsf{Vect}_\R, \R)$: the pairings are bilinear (trace,
$L^2$ inner product, Fourier integral) and the algebra is standard
linear algebra.
At $\varepsilon = 0$, the Maslov limit deforms the ambient category to
$\Chu(\mathsf{Set}, \R_{\max})$: bilinear pairings become tropical
(max-plus), the heat equation becomes Hamilton--Jacobi, and the
Gaussian kernel becomes the min-cost envelope.
The barrier parameter $\varepsilon$ interpolates continuously between
these two categorical regimes---it is the deformation parameter of the
$*$-autonomous structure itself.
\end{remark}

\begin{remark}[Cayley--Dickson tower and the division-algebra hierarchy]
\label{rmk:cayley-dickson}
The progression
$\R \subset \mathbb{C} \subset \mathbb{H} \subset \mathbb{O}$
mirrors the structural layers of the quadruped controller:
\begin{enumerate}[nosep]
\item $\R$ (1-dimensional): the scalar cost $J \in \R$---totally
  ordered, commutative, associative.
\item $\mathbb{C}$ (2-dimensional): the gait phase
  $e^{i\theta} \in S^1$---no preferred direction (loss of ordering),
  the Pontryagin dual $\widehat{G}$ of item~(iii).
\item $\mathbb{H}$ (4-dimensional): the base orientation
  $q \in S^3 \subset \mathbb{H}$, modulated by four feet generating
  torques in $\mathrm{Im}(\mathbb{H}) \cong \R^3$
  (\cref{rem:quaternion-closure}).
  Loss of commutativity: the order of rotations matters.
\item $\mathbb{O}$ (8-dimensional): conjecturally, the eight Chu
  objects of \cref{prop:eight-chu}.
  Loss of associativity: forward composition
  $\beta \to \sigma \to x$ is deterministic (a function),
  but backward decomposition $x(T) \to \beta$ depends on the
  grouping---the variational principle
  \eqref{eq:least-action} selects among multiple factorizations.
\end{enumerate}
At each Cayley--Dickson doubling, a structural property is lost;
at each level of the controller hierarchy, a new form of
non-uniqueness appears.

\medskip\noindent\emph{Cayley--Dickson splitting.}
The Cayley--Dickson construction doubles
$\mathbb{H}$ to $\mathbb{O} = \mathbb{H} \oplus \mathbb{H}\,e_4$
by adjoining a unit $e_4$ that anticommutes with every quaternion
imaginary unit.
This suggests a partition of the eight Chu objects into two
quaternionic halves:
\begin{itemize}[nosep]
\item \emph{Control-loop quaternion} ($\mathbb{H}$ component):
  Pontryagin--Fourier~(iii) as the identity~$1$ (representation
  theory furnishes the scalar frame), together with the three
  imaginary units that form the physical control loop---%
  Barrier--Boltzmann~(vii) as $e_1$ (contact constraints),
  Riccati~(i) as $e_2$ (optimal feedback),
  and Lie--Poisson~(vi) as $e_3$ (momentum dynamics).
  Their internal product $e_1 \cdot e_2 = e_3$ is the contact
  Jacobian triad constructed below.
\item \emph{Spectral-analytic quaternion}
  ($\mathbb{H}\,e_4$ component):
  the four objects providing mathematical infrastructure for the
  control loop---%
  Hopf--Cole~(ii) as the doubling unit~$e_4$
  (the linearization bridge between nonlinear control and
  linear analysis),
  Hardy~(iv) as $e_5$ (causal structure),
  Peter--Weyl~(v) as $e_6$ (spectral decomposition),
  and Action--State~(viii) as $e_7$
  (contact schedule parameterization).
\end{itemize}
The doubling unit $e_4 = \text{Hopf--Cole}$ is the exponential
bridge $V \leftrightarrow \psi = e^{-V/(2\varepsilon)}$:
it converts nonlinear control ($\mathbb{H}$) into linear
analysis ($\mathbb{H}\,e_4$), exactly as the Cayley--Dickson
imaginary unit converts quaternion arithmetic into octonion
arithmetic.

\medskip\noindent\emph{Addition vs.\ multiplication: the physical
root of non-associativity.}
The two halves do not compose associatively because addition and
multiplication are incompatible operations: the superposition of
forces is additive ($f_1 + f_2$), but the composition of the
rotations they induce is multiplicative
($\exp(\tau_1 \Delta t)\,\exp(\tau_2 \Delta t)
\neq \exp((\tau_1 + \tau_2)\Delta t)$ in general).
This is the Baker--Campbell--Hausdorff obstruction:
$\log(e^X e^Y) = X + Y + \tfrac{1}{2}[X,Y] + \cdots$, where the
commutator terms break associativity of the naive
``add-then-exponentiate'' composition.
In the Cayley--Dickson language, the $\mathbb{H}$ and
$\mathbb{H}\,e_4$ halves fail to associate because
$e_4$ anticommutes with the quaternion units---precisely the
algebraic shadow of the BCH remainder.

\medskip\noindent\emph{$\varepsilon$ as the unique
addition--multiplication bridge.}
The nonuple parameter $\varepsilon$ of \S\ref{sec:articulated}
appears in every conversion between the additive and multiplicative
halves:
\begin{itemize}[nosep]
\item Hopf--Cole: $V \leftrightarrow \psi = e^{-V/(2\varepsilon)}$
  (value function $\leftrightarrow$ linear solution);
\item Sigmoid: $\beta \leftrightarrow \sigma =
  (1 + e^{-\beta})^{-1}$
  (log-contact mode $\leftrightarrow$ contact fraction);
\item Log-barrier: $h \leftrightarrow -\varepsilon\log h$
  (constraint slack $\leftrightarrow$ barrier cost);
\item Boltzmann: $J \leftrightarrow w \propto e^{-J/\varepsilon}$
  (cost $\leftrightarrow$ sampling weight).
\end{itemize}
These are four manifestations of a single map
$x \mapsto e^{-x/\varepsilon}$ (or its inverse
$w \mapsto -\varepsilon \log w$).
The parameter $\varepsilon$ is the \emph{exchange rate} between
addition and multiplication---the analogue of $\hbar$ in the
Feynman path integral, where it converts the additive Lagrangian
$\int L\,dt$ to the multiplicative amplitude $e^{iS/\hbar}$.

\medskip\noindent\emph{Explicit assignment and the Fano plane.}
The proposed correspondence between octonion basis elements and
Chu objects is:

\begin{center}
\small
\begin{tabular}{cll}
\hline
Basis & Chu object & Role \\
\hline
$1$ & Pontryagin--Fourier~(iii) & identity / representation frame \\
$e_1$ & Barrier--Boltzmann~(vii) & contact constraints \\
$e_2$ & Riccati~(i) & optimal feedback \\
$e_3$ & Lie--Poisson~(vi) & momentum dynamics \\
$e_4$ & Hopf--Cole~(ii) & linearization bridge (doubling unit) \\
$e_5$ & Hardy~(iv) & causal structure \\
$e_6$ & Peter--Weyl~(v) & spectral decomposition \\
$e_7$ & Action--State~(viii) & contact schedule \\
\hline
\end{tabular}
\end{center}

\noindent
The seven Fano plane triads---cyclic triples $(e_a, e_b, e_c)$
with $e_a \cdot e_b = e_c$---are:

\begin{center}
\footnotesize
\begin{tabular}{clll}
\hline
\# & Product & Physical name & Mediator \\
\hline
T1 & $e_1 {\cdot}\, e_2 {=} e_3$ (BB${\times}$Ri$=$LP)
   & Contact Jacobian & $r_k \times f_k$ \\
T2 & $e_1 {\cdot}\, e_4 {=} e_5$ (BB${\times}$HC$=$Ha)
   & Feynman--Kac proj.\ & Szeg\H{o} $P_+$ \\
T3 & $e_1 {\cdot}\, e_7 {=} e_6$ (BB${\times}$AS$=$PW)
   & Gait spectral analysis & Fourier sign \\
T4 & $e_2 {\cdot}\, e_4 {=} e_6$ (Ri${\times}$HC$=$PW)
   & Riccati spectral dec.\ & symplectic $J$ \\
T5 & $e_2 {\cdot}\, e_5 {=} e_7$ (Ri${\times}$Ha$=$AS)
   & Causal fb.\ $\to$ contacts & spec.\ factor sign \\
T6 & $e_3 {\cdot}\, e_4 {=} e_7$ (LP${\times}$HC$=$AS)
   & Momentum linearization & Lie bracket \\
T7 & $e_3 {\cdot}\, e_6 {=} e_5$ (LP${\times}$PW$=$Ha)
   & Coadjoint orbit caus.\ & conjugation \\
\hline
\end{tabular}
\end{center}

\noindent
Each triad carries a natural antisymmetric bilinear form (cross
product, symplectic form, Lie bracket, Szeg\H{o} projection,
Fourier inversion) that mediates the product; swapping inputs
through this mediator reverses the sign, yielding
anti-commutativity $e_a \cdot e_b = -e_b \cdot e_a$
automatically.
Triad~T1 is the contact Jacobian constructed in detail below;
the remaining six triads are identified by their physical content
but their full Chu tensor-product verification is deferred.

\medskip\noindent\emph{Three doublings, three laws.}
The three Cayley--Dickson doublings correspond to Newton's three
laws of mechanics:
\begin{enumerate}[nosep]
\item $\R \to \mathbb{C}$ (adjoin $e_1 = \text{BB}$):
  \textbf{Law~I} (inertia).
  The contact constraint defines the reference frame; without
  external force, the barrier keeps feet planted.
  Loss of total ordering: multiple equilibria coexist.
\item $\mathbb{C} \to \mathbb{H}$ (adjoin $e_2 = \text{Ri}$,
  derive $e_3 = e_1 \cdot e_2 = \text{LP}$):
  \textbf{Law~II} ($F = dp/dt$).
  Riccati computes optimal force, Lie--Poisson tracks the
  resulting momentum.
  Loss of commutativity: the order of force application matters.
\item $\mathbb{H} \to \mathbb{O}$ (adjoin $e_4 = \text{HC}$,
  derive $e_5, e_6, e_7$):
  \textbf{Law~III} (action $= -$reaction).
  Hopf--Cole linearization enables spectral analysis of the
  control loop; non-associativity appears because multi-body
  contact changes the system topology.
  Loss of associativity: grouping of composed transformations
  matters.
\end{enumerate}

\medskip\noindent\emph{Self-interaction: $e_i^2 = -1$ as
Green's function.}
In the Chu framework, $e_i^2$ is the tensor product of a Chu
object with itself: the object acts simultaneously as source~($A$)
and field~($X$), coupled through the pairing~$r$.
The self-duality isomorphism $\varphi \colon A \xrightarrow{\sim} X$
implicit in each pair carries an antisymmetric structure:
the Riccati Hamiltonian is $J$-symmetric (symplectic),
the Lie--Poisson bracket is skew on $\mathfrak{g}^*$,
and the Szeg\H{o} projection $P_+$ satisfies
$P_+ + P_- = I$ with cross-term $P_+ P_-^* = 0$.
In each case, the sign of the self-pairing is
$r(\varphi(a), a) = -\|a\|^2$---the Green's equation
$L \cdot G = -\delta$ evaluated on the diagonal.
The universal $-1$ is the hallmark of symplectic/antisymmetric
self-duality, not an algebraic accident.

\medskip
Hurwitz's theorem ($\mathbb{O}$ is the last normed division
algebra) would, if the correspondence is exact, imply that eight
Chu dualities are \emph{maximal}---no ninth independent duality
exists without the algebraic structure collapsing.
Wu et al.~\cite{WuOctonion} provide empirical evidence that
octonion-structured cross-correlations (``multi-task learning via
the octonion multiplication table'') outperform quaternion or
complex alternatives---consistent with the hypothesis that eight
coupled channels carry more structure than four.

\medskip\noindent\emph{Example: the contact Jacobian as a Chu triad.}
As evidence for the conjectured multiplication, consider the
triple (Barrier--Boltzmann~(vii), Riccati~(i), Lie--Poisson~(vi)).
In the centroidal model (\S\ref{sec:articulated}), each stance
foot~$k$ contributes a force $\sigma_k f_k$ to the body through
the moment arm $r_k$; the net generalized force is
$\mu = \sum_k \sigma_k\, r_k \times f_k \in \mathfrak{g}^*$.
Define:
\begin{itemize}[nosep]
\item \emph{Forward map} $\phi$: given contact modes
  $\sigma \in (0,1)^4$ (from Barrier--Boltzmann) and optimal
  forces $f_k = -K_k(t)\,\delta x$ (from Riccati), produce the
  body torque $\mu = \sum_k \sigma_k\, r_k \times f_k
  \in \mathfrak{g}^*$ (in Lie--Poisson).
\item \emph{Backward map} $\phi^*$: given a required body torque
  $\mu \in \mathfrak{g}^*$, decompose it into per-foot
  contributions $\sigma_k\, r_k \times f_k$ weighted by
  $\sigma_k$---the force-distribution problem.
\end{itemize}
The adjunction condition
$\langle \mu, \xi \rangle_{\mathfrak{g}}
= \sum_k \sigma_k\, \langle r_k \times f_k, \xi \rangle$
is d'Alembert's principle of virtual work: the body absorbs
exactly the power delivered by the contact-gated feet.
The bilinearity in $(\sigma_k, f_k)$---contact mode times
control force---is the tensor-product structure:
neither factor alone determines the body torque.

\medskip\noindent\emph{Analytic criterion: Jensen's gap as a
non-associativity measure.}
The strictly convex barrier
$V(h) = -\varepsilon \log h$ satisfies Jensen's inequality:
$\mathbb{E}[V(h)] \geq V(\mathbb{E}[h])$, with equality if and
only if $h$ is constant.
Consider the associator for the triple
$(e_1, e_2, e_5) = (\text{BB}, \text{Ri}, \text{Ha})$:
\begin{itemize}[nosep]
\item \emph{Left bracket}
  $(e_1 \cdot e_2) \cdot e_5 = e_3 \cdot e_5 = -e_6$:
  solve the barrier-gated Riccati feedback (producing
  Lie--Poisson dynamics $e_3$), \emph{then} project onto the
  causal Hardy component---yielding
  $\mathbb{E}[V_{\mathrm{barrier}}(\cdot)]$.
  The result is $-e_6 = -\text{PW}$ (negative spectral content).
\item \emph{Right bracket}
  $e_1 \cdot (e_2 \cdot e_5) = e_1 \cdot e_7 = +e_6$:
  first compose Riccati with Hardy splitting (yielding
  Action--State contact selection $e_7$), \emph{then} apply the
  barrier---yielding
  $V_{\mathrm{barrier}}(\mathbb{E}[\cdot])$.
  The result is $+e_6 = +\text{PW}$ (positive spectral content).
\end{itemize}
The associator is $-2e_6$: the discrepancy between the two
evaluation orders is measured by Peter--Weyl spectral content---the
frequency-domain representation of the gait.
The difference is the Jensen gap:
\begin{equation}\label{eq:jensen-gap}
  \Delta_J
  \;=\; \mathbb{E}[-\varepsilon \log h]
  - (-\varepsilon \log \mathbb{E}[h])
  \;=\; \varepsilon \log
  \frac{\overline{h}}{\exp(\overline{\log h})}
  \;\geq\; 0,
\end{equation}
where $\overline{h}$ is the arithmetic mean and
$\exp(\overline{\log h})$ is the geometric mean of the contact
schedule over one stride.
This ratio---arithmetic mean over geometric mean---is the
signature of the Szeg\H{o} structure from Part~I:
the geometric mean $G(f) = \exp(\int_\bT \log f\, d\theta)$
of \cref{def:cycle} appears in the denominator, and
$\Delta_J = 0$ if and only if the contact schedule is
constant, i.e., the Gaussian fixed point of
\cref{thm:holonomy-trichotomy}(i) where $H(f) = 0$.
For any non-trivial periodic gait ($H(f) > 0$), the Jensen gap
is strictly positive: the two evaluation orders disagree, and
the Chu tensor product is not associative.

We do not claim that $\Delta_J$ equals the full octonion
associator; the holonomy $H(f)$ of \cref{def:cycle},
which involves all weighted Fourier coefficients
$\sum_{k \geq 1} k\,|c_k|^2$, is a finer invariant than the
scalar AM/GM ratio.
But both vanish on exactly the same set ($f = \mathrm{const}$),
and the Jensen gap provides an analytic certificate:
\emph{non-trivial locomotion ($H(f) > 0$) is a necessary
condition for the eight Chu dualities to exhibit
non-associative coupling}.

\medskip
We state the correspondence as a conjecture with the following
supporting evidence:
\begin{enumerate}[label=(\alph*),nosep]
\item An explicit assignment
  (PF~$\mapsto 1$,
  BB~$\mapsto e_1$,
  Ri~$\mapsto e_2$,
  LP~$\mapsto e_3$,
  HC~$\mapsto e_4$,
  Ha~$\mapsto e_5$,
  PW~$\mapsto e_6$,
  AS~$\mapsto e_7$)
  that respects the Cayley--Dickson splitting
  $\mathbb{O} = \mathbb{H} \oplus \mathbb{H}\,e_4$
  (control-loop vs.\ spectral-analytic quaternions).
\item One triad fully constructed: T1 (contact Jacobian,
  $e_1 \cdot e_2 = e_3$) with explicit forward/backward maps
  and the d'Alembert adjunction condition.
\item All seven triads identified with physical content and
  antisymmetric mediators (table above).
\item The self-product $e_i^2 = -1$ follows from
  symplectic self-duality (Green's function sign).
\item Anti-commutativity $e_a \cdot e_b = -e_b \cdot e_a$
  follows from the antisymmetric mediator in each triad.
\item The Jensen gap~\eqref{eq:jensen-gap} certifies that
  non-trivial locomotion ($H(f) > 0$) implies
  non-associativity; the concrete associator
  $-2e_6$ for the triple $(e_1, e_2, e_5)$
  identifies Peter--Weyl spectral content as the obstruction.
\end{enumerate}
What remains open is the rigorous verification that each of the
seven physical morphisms satisfies the universal property of the
Chu tensor product---specifically, that the forward and backward
maps of each triad compose to yield the correct pairing on the
third Chu object.
\end{remark}
